% Generated draft: requires mathematical and source review.
\documentclass[11pt]{article}
\usepackage[margin=25mm]{geometry}
\usepackage{fontspec}
\setmainfont{Latin Modern Roman}
\usepackage{amsmath,amssymb,mathtools,unicode-math}
\setmathfont{Latin Modern Math}
\usepackage{enumitem,needspace,xurl,hyperref}
\setlength{\parindent}{0pt}
\setlength{\parskip}{5pt}
\setlength{\emergencystretch}{3em}
\newcommand{\N}{\mathbb{N}}
\newcommand{\Z}{\mathbb{Z}}
\newcommand{\Q}{\mathbb{Q}}
\newcommand{\R}{\mathbb{R}}
\newcommand{\C}{\mathbb{C}}
\newcommand{\F}{\mathbb{F}}
\newcommand{\A}{\mathbb{A}}
\newcommand{\PP}{\mathbb{P}}
\newcommand{\E}{\mathbb{E}}
\newcommand{\eps}{\varepsilon}
\newcommand{\dd}{\,\mathrm{d}}
\begin{document}
\section*{UnsolvedMath ID: 30000636}
\textbf{Source problem code:} catalogue entry
\subsection{Definitions and mathematical statement}
The source asks for the best H\"older exponents in stability estimates for recovery of convex bodies from tomographic data formed from sections, projections and averages of those data. It separately asks for the optimal dimensional dependence, conjectured to be linear in \(1/n\) rather than quadratic. No particular tomography operator or stability inequality is reproduced.

A precise instance requires an admissible class \(\mathcal K_n\) of convex bodies in \(\R^n\), a specified data map \(T_n\), a body distance \(d_{\mathcal K}\) and a data distance \(d_{\mathcal D}\). The desired estimate would have the form
\[
 d_{\mathcal K}(K,L)\leq C\,
 d_{\mathcal D}(T_nK,T_nL)^{\alpha_n}
 \quad(K,L\in\mathcal K_n),
\]
possibly only when the data distance is below a specified tolerance. The best exponent can then be defined as the supremum of \(\alpha>0\) for which such a bound holds with one allowed constant \(C\), uniform over the admissible pairs and the prescribed error range. The class, constant dependencies and local-versus-global regime must remain fixed while optimizing the exponent. Sharpness requires a matching obstruction to every larger exponent, not just proof of one bound.

Sections and projections need their dimensions, directions and offsets, while averages need their averaging measure and retained parameters. Bounds for these different data maps cannot be conflated or treated as bounds for their combined data without specification. Hausdorff, symmetric-difference and other distances are also distinct choices. Centering, normalization, geometric size bounds and any symmetry or regularity assumptions are absent. If the data do not distinguish two bodies in the chosen class, an estimate forcing zero body distance for equal data cannot hold unless the class or equivalence is adjusted; no injectivity theorem is presumed here.

The dimensional question concerns \(\alpha_n\), not simply growth of the prefactor \(C\). The phrase linear in \(1/n\) rather than quadratic suggests distinguishing a scale like \(1/n\) from one like \(1/n^2\), but supplies neither exact formulas nor uniform comparison constants or a specified operator to which that conjecture applies. Those are not invented. Both the general sharp-exponent task and this particular dimensional conjecture are retained.

The literature triage checked 2026-08-21 reports no authoritative complete resolution of the extracted claim located in its search. The open label and verification note describing typography/notation cleanup remain dated source metadata, not independently checked current results.
\subsection{Short English statement}
Find sharp powers relating reconstruction error to tomographic-data error, separately for the intended section, projection and averaged-data problems. Then determine how those powers depend on dimension, including the proposed \(1/n\) rather than \(1/n^2\) scale. The operators, metrics and admissible body classes are missing.
\subsection{Sources}
\begin{itemize}
\item Local UnsolvedMath registry, permanent ID 30000636, source code catalogue entry, ``Optimal Stability Exponents in Convex-Body Tomography,'' including the general sharp-exponent task, dimensional conjecture and verification metadata.
\item \emph{Konvexgeometrie (2007)}, Oberwolfach Reports, pp.\ 3321--3396, DOI 10.4171/owr/2006/56: \url{https://doi.org/10.4171/owr/2006/56}. This is also the sole supplied literature reference. The citation uses 2007, while the DOI and proposed-year metadata use 2006.
\item Literature triage checked 2026-08-21; record updated 2026-09-16. No external source has been visited or independently checked here.
\end{itemize}
\end{document}
